\chapter{Complexity — Diagonalization}
%%% generated by the blueprint pipeline from TCSlib.Complexity.Uncomputability.Diagonalization
%%%%%%%%%%%%%%%%%%%%%%%%%%%%%%%%%%%%%%%%%%%%%%%%%%%%%%%%%%%%%%%%%%%%%%%%%%%%%%%
\section{Overview}

This module proves uncomputability by diagonalization [AB09, \S 1.5,
Theorem 1.10]: it defines the diagonal function $\mathrm{UC}$ of a
representation scheme for Turing machines and shows that $\mathrm{UC}$ is not
computable, for every such scheme --- effective or not. The result uses only
machines-as-strings, not the universal machine.

%%%%%%%%%%%%%%%%%%%%%%%%%%%%%%%%%%%%%%%%%%%%%%%%%%%%%%%%%%%%%%%%%%%%%%%%%%%%%%%
\section{Declarations}

\begin{definition}[Diagonal function of a representation scheme]
\lean{Complexity.UC}
\label{Complexity.UC}
\leanok
\uses{Turing.CodeTM.toFinTM, Turing.FinTM.ComputesInTime, Turing.MachineCode}
Fix a representation scheme $c$ for Turing machines in code normal form (one
work tape, binary alphabet): a total decoding assigning to every bit string a
machine, together with an encoding that the decoding inverts. The diagonal
function $\mathrm{UC}_c : \{0,1\}^* \to \{0,1\}$ maps a bit string $\alpha$ to
$\mathrm{false}$ if the machine that $\alpha$ denotes under $c$, run on the
input $\alpha$ itself, halts after some finite number of steps with completed
output the one-bit string $[\mathrm{true}]$ --- that is, if the machine accepts
its own representation. In every other case --- the machine diverges on
$\alpha$, or halts with any completed output other than $[\mathrm{true}]$ ---
the value is $\mathrm{true}$. ([AB09, \S 1.5] writes: $UC(\alpha) = 0$ if
$M_\alpha(\alpha) = 1$, and $UC(\alpha) = 1$ otherwise.)
\end{definition}

\begin{theorem}[Characterization of the value false of the diagonal function]
\lean{Complexity.UC_eq_false_iff}
\label{Complexity.UC_eq_false_iff}
\leanok
\uses{Complexity.UC, Turing.CodeTM.toFinTM, Turing.FinTM.ComputesInTime, Turing.MachineCode}
\difficulty{2}
Let $c$ be a representation scheme for Turing machines in code normal form (one
work tape, binary alphabet; a total decoding of bit strings to machines,
inverted by the encoding), and let $\alpha$ be a bit string. Then the diagonal
function of $c$ takes the value $\mathrm{false}$ at $\alpha$ --- i.e.\
$\mathrm{UC}_c(\alpha) = \mathrm{false}$ --- if and only if there exists a
number of steps $t$ within which the machine denoted by $\alpha$ under $c$,
run on the input $\alpha$ itself, halts with completed output the one-bit
string $[\mathrm{true}]$.
\end{theorem}

\begin{theorem}[Characterization of the value true of the diagonal function]
\lean{Complexity.UC_eq_true_iff}
\label{Complexity.UC_eq_true_iff}
\leanok
\uses{Complexity.UC, Turing.CodeTM.toFinTM, Turing.FinTM.ComputesInTime, Turing.MachineCode}
\difficulty{2}
Let $c$ be a representation scheme for Turing machines in code normal form (one
work tape, binary alphabet; a total decoding of bit strings to machines,
inverted by the encoding), and let $\alpha$ be a bit string. Then the diagonal
function of $c$ takes the value $\mathrm{true}$ at $\alpha$ --- i.e.\
$\mathrm{UC}_c(\alpha) = \mathrm{true}$ --- if and only if there is no number
of steps $t$ within which the machine denoted by $\alpha$ under $c$, run on the
input $\alpha$ itself, halts with completed output the one-bit string
$[\mathrm{true}]$: the machine diverges on $\alpha$, or halts with a completed
output other than $[\mathrm{true}]$.
\end{theorem}

\begin{theorem}[Uncomputability of the diagonal function]
\lean{Complexity.UC_not_computable}
\label{Complexity.UC_not_computable}
\leanok
\uses{Complexity.Computable, Complexity.UC, Complexity.UC_eq_false_iff, Turing.CodeTM.toFinTM, Turing.FinTM.Computes.exists_computesFunInTime, Turing.FinTM.ComputesInTime, Turing.FinTM.ComputesInTime.output_unique, Turing.MachineCode, Turing.MachineCode.decode_encode, Turing.exists_codeTM}
\difficulty{5}
Let $c$ be any representation scheme for Turing machines in code normal form
(one work tape, binary alphabet; a total decoding of bit strings to machines,
inverted by the encoding), effective or not, and let $\mathrm{UC}_c$ be its
diagonal function: $\mathrm{UC}_c(\alpha) = \mathrm{false}$ if the machine
denoted by $\alpha$ under $c$ halts on the input $\alpha$ itself with completed
output the one-bit string $[\mathrm{true}]$, and
$\mathrm{UC}_c(\alpha) = \mathrm{true}$ otherwise. Then the total function on
bit strings $\alpha \mapsto [\mathrm{UC}_c(\alpha)]$, returning the one-bit
string whose single bit is the value of the diagonal function, is not
computable: no finite Turing machine over the binary alphabet halts on every
input $\alpha$ with $[\mathrm{UC}_c(\alpha)]$ on its output tape [AB09,
Theorem 1.10].
\end{theorem}
